\section{User study}\label{user_study}
A user study is an approach taken to understand the requirements of users in relation to your application \parencite{user_study_definition}. It can be used to better understand user requirements for those who will be using the system and consequently improve design decisions.

\subsection{Study design}
The user study was designed with the goal of evaluating the effectiveness of the tool, \lintex, and identifying areas of improvement for the tool. To meet these goals, the user study was split into six key stages:
\begin{enumerate}
	\item An introductory stage where the goal is to identify where the participant's knowledge lies in related topics.\label{user_study:intro_questions}
	\item An initial listening of the \gls{inaccessible_doc}, which was produced for this user study.\label{user_study:initial_listening}
	\item A questioning stage based on the initial listening of the \gls{inaccessible_doc}.\label{user_study:initial_questions}
	\item A think-aloud session for understanding the accessibility of \lintex.\label{user_study:using_lintex}
	\item A final listening of the newly created \gls{accessible_doc} by using LinTeX.\label{user_study:final_listening}
	\item A questioning stage based on the final listening of the \gls{accessible_doc}.\label{user_study:final_questions}
\end{enumerate}
These stages allowed for a comparison of the before and after of using \lintex, through stages \ref{user_study:initial_listening} and \ref{user_study:final_listening}, as well as associated questioning stages \ref{user_study:initial_questions} and \ref{user_study:final_questions}, allowing both qualitative and quantitative data to be obtained. Additionally, the think-aloud session provides an opportunity to judge the usability of \lintex.

Finer details for the user study can be found in the \fullref{appendix:user_study_script}, with the reasoning behind their choices found in the \fullref{backg:user_study_design}.
\subsection{Hypothesis}
There were several hypotheses that this user study aimed to prove or disprove:
\begin{itemize}
	\item[\textbf{H1}] The \gls{accessible_doc} created by \lintexspaced showed an improvement from the \gls{inaccessible_doc} initially used in terms of the document's accessibility for screen reader users.\label{user_study:hypothesis:improved_doc_accessibility}
	\item[\textbf{H2}] \lintexspaced is largely accessible.\label{user_study:hypothesis:lintex_accessible}
	\item[\textbf{H3}] \lintexspaced provided useful error messages.\label{user_study:hypothesis:errors}
	\item[\textbf{H4}] The \gls{inaccessible_doc} used in the initial listening is not memorable.\label{user_study:hypothesis:memorable}
\end{itemize}

\subsection{Results}


\subsection{Analysis}
\subsection{Hypothesis testing}